








































\documentclass[11pt]{article}

\usepackage{amsmath,amssymb,amsthm}
\usepackage{geometry}
\usepackage{hyperref}

\geometry{margin=2.5cm}

\title{
Lois expérimentales du PCH :
Analyse globale des régimes dimension--calibration,
stabilité géométrique et applications aux LLM
}

\author{}
\date{}

\begin{document}

\maketitle


\begin{abstract}

Cette étude synthétise l'ensemble des expériences PCH historiques
et des nouvelles expériences ex14--ex20.

Les résultats révèlent que PCH n'est pas simplement un mécanisme
d'amélioration de représentation mais un opérateur géométrique
contextuel qui transforme un espace latent selon une contrainte
d'information.

La variable déterminante n'est pas uniquement la dimension latente
mais le rapport :

\[
\kappa =
\frac{dim}{n_{\mathrm{calib\_per\_context}}}
\]

L'étude montre l'existence d'un régime critique lorsque :

\[
dim=n_{\mathrm{calib\_per\_context}}
\]

où apparaissent une stabilité maximale, une projection quasi exacte
et une forte amélioration oracle.

\end{abstract}


\section{Modèle conceptuel}

Le système applique une transformation :

\[
z \rightarrow P_c(z)
\]

où :

\[
c
\]

représente un contexte latent.

Le projecteur dépend donc du contexte :

\[
P_c:\mathbb{R}^{D}\rightarrow\mathbb{R}^{D}
\]

mais modifie la géométrie interne des représentations.


Contrairement à une simple transformation statistique,
PCH agit comme une contraction vers une variété contextuelle.


\section{Loi 1 : attracteur universel de normalisation}


Toutes les expériences convergent vers :

\[
\gamma=1
\]

et :

\[
\tau=3.4
\]


Les valeurs restent invariantes malgré :

\begin{itemize}
\item changement de ridge ;
\item suppression du régularisateur ;
\item nombre d'itérations ;
\item changement de dimension.
\end{itemize}


La conséquence est :

\[
\boxed{
La normalisation impose une géométrie cible universelle.
}
\]


Le système ne cherche pas une optimisation classique.
Il force une configuration géométrique stable.


\section{Loi 2 : contraction et stabilité dynamique}


Toutes les expériences donnent :

\[
N_{\mathrm{violations}}=0
\]


pour la monotonie énergétique.


Donc :

\[
E(P_c(x))\leq E(x)
\]


Les résidus point fixe sont :

\[
10^{-17}
\]

à :

\[
10^{-12}
\]


même dans les régimes difficiles.


Ainsi :

\[
\boxed{
PCH\ est un opérateur contractant.
}
\]


La stabilité ne dépend pas seulement de la dimension.


Elle dépend de la cohérence entre :

\[
dim
\]

et :

\[
n_{\mathrm{calib}}
\]


\section{Loi 3 : compression géométrique}


Les expériences montrent une réduction massive :

\[
\Delta d =
d_{avant}-d_{après}
\]


avec :

\[
\frac{\Delta d}{d_{avant}}
\approx 90\%-95\%
\]


Exemples :


\[
0.036\rightarrow0.0028
\]


(ex14)


et :

\[
0.232\rightarrow0.027
\]


(ex18)


Donc :

\[
\boxed{
PCH rapproche les représentations d'une variété contextuelle.
}
\]


\section{Découverte principale :
loi dimension--calibration}


La variable fondamentale est :

\[
\kappa=
\frac{dim}{n_{\mathrm{calib\_per\_context}}}
\]


Les expériences révèlent trois régimes.


\subsection{Régime équilibré}


Lorsque :

\[
dim=n_{\mathrm{calib}}
\]


on observe :

\[
\kappa=1
\]


Cas :

\[
(4,4)
\]

(ex14)


Résultats :

\[
\rho=0.2272727
\]


\[
gap_{operator}\approx10^{-12}
\]


\[
gap_{projector}\approx10^{-12}
\]


Le projecteur devient presque idéal.


Conséquence :

\[
\boxed{
dim=n_{\mathrm{calib}}
\Rightarrow
géométrie optimale
}
\]


Il existe un alignement parfait entre :

\[
\text{nombre de degrés de liberté}
\]

et :

\[
\text{nombre de contraintes contextuelles}.
\]


\section{Régime sous-dimensionné}


Lorsque :

\[
dim<n_{\mathrm{calib}}
\]


Exemple :

\[
dim=1,n_{\mathrm{calib}}=1
\]

(ex18)


donne :

\[
\rho=0.089
\]


La stabilité reste parfaite :

\[
fixed\ point\approx10^{-17}
\]


mais la capacité représentative chute.


Le projecteur conserve trop peu de directions.


Le résultat :

\[
accuracy_{oracle}=1
\]


mais :

\[
accuracy_{routed}<baseline
\]


Donc :

\[
\boxed{
La projection est correcte mais la représentation est trop compressée.
}
\]


\section{Régime surdimensionné}


Lorsque :

\[
dim>n_{\mathrm{calib}}
\]


Exemple :

\[
dim=32,n_{\mathrm{calib}}=48
\]


Même si :

\[
\kappa<1
\]


la calibration devient suffisante.


On observe :

\[
operator\ gap
\approx10^{-4}
\]


au lieu de :

\[
10^{-12}
\]


Le système reste stable mais la projection devient moins exacte.


\section{Nouvelle loi générale}


Les résultats imposent :

\[
\boxed{
Qualité(PCH)=f
\left(
\frac{dim}{n_{\mathrm{calib}}}
\right)
}
\]


La zone optimale est :

\[
\boxed{
\frac{dim}{n_{\mathrm{calib}}}\approx1
}
\]


Trop faible :

\[
\frac{dim}{n_{\mathrm{calib}}}\ll1
\]

perte d'information.


Trop élevé :

\[
\frac{dim}{n_{\mathrm{calib}}}\gg1
\]

sous-détermination.


\section{Expérience ex20 : saturation par sur-calibration}


Cas :

\[
dim=1
\]

mais :

\[
n_{\mathrm{calib}}=120
\]


Le système atteint :

\[
\rho=0.22727
\]


mais :

\[
distance_{avant}
\approx0
\]


Donc :

\[
PCH
\]

n'a plus rien à corriger.


La transformation devient quasi identité.


Conclusion :


\[
\boxed{
Une calibration excessive peut annuler le gain.
}
\]


La calibration doit correspondre à la complexité réelle.


\section{Oracle versus routage}


Toutes les expériences montrent :


\[
accuracy_{oracle}
>
accuracy_{routed}
\]


Donc :

le problème principal n'est pas :

\[
P_c
\]


mais :

\[
\hat c
\]


l'estimation du contexte.


L'amélioration future doit porter sur :

\[
x\rightarrow c
\]

plutôt que :

\[
P_c
\]


\section{Conséquences pour les LLM}


Un LLM classique :

\[
z=f(x)
\]


devient :

\[
z'=P_c(f(x))
\]


où :

\[
c
\]

est un contexte latent.


\subsection{Adaptation sans fine tuning}


PCH peut remplacer certains adaptateurs.


Au lieu de :

\[
LLM+\mathrm{LoRA}
\]


on obtient :

\[
LLM+P_c
\]


avec plusieurs géométries :

\[
P_{medical}
\]

\[
P_{legal}
\]

\[
P_{scientific}
\]


\subsection{Mémoire contextuelle}


La mémoire devient :

\[
M=
\{P_c(z)\}
\]


Les souvenirs sont stockés
dans une géométrie cohérente.


\subsection{RAG géométrique}


Le pipeline devient :

\[
document
\rightarrow embedding
\rightarrow P_c(embedding)
\]


Le contexte ne sélectionne plus seulement des documents.

Il transforme l'espace de recherche.


\subsection{Détection d'anomalies}


Le score :

\[
A(x)=||x-P_c(x)||
\]


devient un indicateur de :

\[
\text{incohérence contextuelle}
\]


\subsection{Routage multi-experts}


Chaque expert peut avoir :

\[
P_i
\]


Le contexte choisit :

\[
P_{\hat c}
\]


sans changer le modèle principal.


\section{Lois finales révélées}


Les expériences supportent les lois suivantes :


\[
\boxed{
L_1:
\text{La normalisation crée un attracteur géométrique universel}
}
\]


\[
\boxed{
L_2:
PCH est un opérateur contractant stable
}
\]


\[
\boxed{
L_3:
La variable fondamentale est }
\frac{dim}{n_{\mathrm{calib}}}
\]


\[
\boxed{
L_4:
dim=n_{\mathrm{calib}}
\Rightarrow
projection optimale
}
\]


\[
\boxed{
L_5:
Le routage du contexte est le vrai goulot d'étranglement
}
\]


\section{Conclusion}


Les expériences montrent que PCH introduit une nouvelle manière
de contrôler les représentations neuronales.

La dimension seule n'est pas le facteur déterminant.

Le facteur critique est l'accord entre :

\[
\text{capacité géométrique}
\]

et :

\[
\text{information de calibration}.
\]


Cette loi ouvre une voie pour des LLM capables de changer
leur géométrie interne dynamiquement,
sans modification des poids du modèle.


\end{document}